% ============================================================
%  main.tex — Documento principal de la tesis
%  Compilación: XeLaTeX + biber   (estilo de citas: IEEE)
%
%  Comandos útiles:
%    make          → compila (genera build/main.pdf)
%    make watch    → compila en cada cambio
%    make clean    → limpia temporales
% ============================================================

% --- Directivas para editores ---
% !TeX program = xelatex
% !TeX encoding = UTF-8
% !TEX root = main.tex

\documentclass[12pt,a4paper,twoside,openright]{book}

% --- Paquetes, configuración y macros ---
% ============================================================
%  paquetes.tex — Paquetes utilizados en la tesis
%  Compilación: XeLaTeX + biber (citas IEEE)
% ============================================================

% ----- Idioma -----
\usepackage[spanish,es-tabla]{babel}
\usepackage{csquotes}              % comillas tipográficas (requerido por biblatex)

% ----- Fuentes (XeLaTeX / fontspec) -----
% TeX Gyre Termes es un clon de Times, estándar en TeX Live y MiKTeX.
% Para usar "Times New Roman" sustituye por: \setmainfont{Times New Roman}
\usepackage{fontspec}
\setmainfont{TeX Gyre Termes}      % texto (Times)
\setsansfont{TeX Gyre Heros}       % sans
\setmonofont{TeX Gyre Cursor}      % monoespaciada

% ----- Matemáticas -----
\usepackage{amsmath}
\usepackage{amssymb}
\usepackage{mathtools}
\usepackage{unicode-math}
\setmathfont{TeX Gyre Termes Math} % fuentes de matemáticas acordes a Termes

% ----- Figuras, tablas, unidades -----
\usepackage{graphicx}
\usepackage{float}
\usepackage{longtable}
\usepackage{booktabs}
\usepackage{multirow}
\usepackage{array}
\usepackage{caption}
\usepackage{subcaption}

% ----- Diseño de página -----
\usepackage[a4paper,top=3cm,bottom=3cm,left=3.5cm,right=2.5cm,headheight=14pt]{geometry}
\usepackage{setspace}              % interlineado
\usepackage{fancyhdr}              % encabezados / pies
\usepackage{enumitem}              % listas personalizables
\usepackage{etoolbox}              % adaptación de entradas del índice

% ----- Bibliografía (IEEE + biber) -----
% sorting=none → numeración por orden de cita (típico de IEEE)
\usepackage[style=ieee,backend=biber,sorting=none]{biblatex}

% Adaptar book antes de que hyperref envuelva el comando de capítulo.
% Tabla de contenido con una cabecera por capítulo y páginas automáticas.
\addto\captionsspanish{\renewcommand{\contentsname}{TABLA DE CONTENIDO}}

\makeatletter
\newcommand{\l@chaptergroup}[2]{%
  \addvspace{1em}%
  \@dottedtocline{0}{0em}{0em}{\bfseries #1}{\bfseries #2}}
\newcommand{\toclevel@chaptergroup}{0}
\renewcommand{\l@chapter}[2]{%
  \@dottedtocline{0}{0em}{1.8em}{\bfseries #1}{\bfseries #2}}

% Conserva la numeración normal de book y sus enlaces; solo añade la
% cabecera CAPÍTULO N. antes del título numerado en la tabla de contenido.
\patchcmd{\@chapter}
  {\addcontentsline{toc}{chapter}{\protect\numberline{\thechapter}#1}}
  {\addcontentsline{toc}{chaptergroup}{\protect\MakeUppercase{\@chapapp}\ \thechapter.}%
   \addcontentsline{toc}{chapter}{\protect\numberline{\thechapter}#1}}
  {}
  {\PackageError{tesis}{No se pudo configurar la tabla de contenido}%
    {Comprueba la definicion de chapter en la clase book.}}
\makeatother


% ----- Enlaces / metadatos del PDF (debe ir al final) -----
\usepackage[colorlinks=true,linkcolor=blue,citecolor=blue,urlcolor=blue]{hyperref}

% ============================================================
%  configuracion.tex — Configuración global del documento
% ============================================================

% ----- Nombres en español -----
\renewcommand{\chaptername}{Capítulo}
\renewcommand{\appendixname}{Anexo}

% ----- Profundidad de numeración / índice -----
\setcounter{tocdepth}{3}   % hasta subsubsección
\setcounter{secnumdepth}{3}

% ----- Intelineado -----
\onehalfspacing

% ----- Encabezados y pies de página -----
\pagestyle{fancy}
\fancyhf{}
\fancyhead[LE,RO]{\small\leftmark}   % capítulo en páginas pares/impares
\fancyhead[RE,LO]{\small\rightmark}  % sección
\fancyfoot[C]{\thepage}
\renewcommand{\headrulewidth}{0.4pt}
\renewcommand{\footrulewidth}{0pt}

% ----- Tablas -----
\renewcommand{\arraystretch}{1.3}
\captionsetup{labelfont=bf}

% ----- Espaciado entre párrafos -----
\setlength{\parskip}{0.5em}

% ----- Numeración de ecuaciones, figuras y tablas por capítulo -----
\numberwithin{equation}{chapter}
\numberwithin{figure}{chapter}
\numberwithin{table}{chapter}

% ============================================================
%  macros.tex — Comandos personalizados de la tesis
%  Añade aquí tus propios atajos para mantener el texto limpio.
% ============================================================

% ----- Énfasis / términos técnicos -----
\newcommand{\tema}[1]{\emph{#1}}
\newcommand{\codigo}[1]{\texttt{#1}}
\newcommand{\comando}[1]{\texttt{#1}}

% ----- Comillas tipográficas -----
\newcommand{\en}[1]{``#1''}

% ----- Referencias cruzadas abreviadas -----
\newcommand{\figura}[1]{Figura~\ref{#1}}
\newcommand{\tabla}[1]{Tabla~\ref{#1}}
\newcommand{\ecuacion}[1]{Ecuación~\eqref{#1}}
\newcommand{\capitulo}[1]{Capítulo~\ref{#1}}

% ----- Unidades (si no usas siunitx) -----
\newcommand{\celsius}{\ensuremath{^\circ\text{C}}}
\newcommand{\porciento}{\ensuremath{\%}}

% ----- Valores con unidades rápidos -----
\newcommand{\unidad}[2]{#1~\mathrm{#2}}


% --- Recursos de bibliografía ---
\addbibresource{bibliografia/referencias.bib}

\begin{document}

% ============================================================
%  PARTE PRELIMINAR (portada, dedicatoria, resúmenes, índices)
% ============================================================
\frontmatter

% --- Metadatos de la tesis (título, autor, tutor, palabras clave…) ---
% ============================================================
%  datos_tesis.tex — Metadatos de la tesis  (✏️ EDITA AQUÍ)
%  Estos comandos se usan en la portada, resumen, etc.
% ============================================================

% ----- Título -----
\newcommand{\titulo}{Título de la tesis}
\newcommand{\tituloEn}{Thesis title in English}

% ----- Autor -----
\newcommand{\autor}{Nombre y apellidos del autor}

% ----- Dirección / tutoría -----
\newcommand{\director}{Nombre del director o tutor}
\newcommand{\codirector}{Nombre del co-director (opcional)}

% ----- Institución -----
\newcommand{\universidad}{Nombre de la universidad}
\newcommand{\facultad}{Facultad de ...}
\newcommand{\programa}{Programa académico / carrera}

% ----- Ciudad y fecha -----
\newcommand{\ciudad}{Popayán}
\newcommand{\pais}{Colombia}
\newcommand{\anio}{2026}

% ----- Palabras clave (resumen) -----
\newcommand{\palabrasclave}{palabra clave 1; palabra clave 2; palabra clave 3}
\newcommand{\keywords}{keyword 1; keyword 2; keyword 3}


% ============================================================
%  portada.tex — Carátula / portada
% ============================================================

\begin{titlepage}
    \centering
    \thispagestyle{empty}

    \vspace*{1.5cm}
    {\Large\bfseries \universidad\par}
    \vspace{0.4cm}
    {\large \facultad\par}
    {\large \programa\par}

    \vspace{3.0cm}
    {\LARGE\bfseries \titulo\par}
    \vspace{1.5cm}

    \vfill
    {\large Presentado por:\par}
    {\large\bfseries \autor\par}
    \vspace{0.8cm}

    {\large Dirigido por:\par}
    {\large\bfseries \director\par}

    \vfill
    {\large \ciudad, \pais\par}
    {\large \anio\par}
\end{titlepage}

% ============================================================
%  dedicatoria.tex — Dedicatoria
% ============================================================

\chapter*{Dedicatoria}
\addcontentsline{toc}{chapter}{Dedicatoria}

\begin{flushright}
    \emph{A \ldots}\\[1cm]
\end{flushright}

% ============================================================
%  agradecimientos.tex — Agradecimientos
% ============================================================

\chapter*{Agradecimientos}
\addcontentsline{toc}{chapter}{Agradecimientos}

\textit{[Escribe aquí tus agradecimientos. Esta sección es opcional.]}

% ============================================================
%  resumen.tex — Resumen en español
% ============================================================

\chapter*{Resumen}
\addcontentsline{toc}{chapter}{Resumen}

\textit{[Escribe aquí el resumen en español, de aproximadamente 200–300 palabras. Debe sintetizar el problema, los objetivos, la metodología, los resultados y las conclusiones.]}

\vspace{0.5cm}
\noindent\textbf{Palabras clave:} \palabrasclave

% ============================================================
%  abstract.tex — Resumen en inglés
% ============================================================

\chapter*{Abstract}
\addcontentsline{toc}{chapter}{Abstract}

\textit{[Write the English abstract here, about 200–300 words.]}

\vspace{0.5cm}
\noindent\textbf{Keywords:} \keywords

% ============================================================
%  lista_abreviaturas.tex — Lista de abreviaturas y siglas
%  Añade aquí tantas filas como necesites.
% ============================================================

\chapter*{Lista de abreviaturas}
\addcontentsline{toc}{chapter}{Lista de abreviaturas}

\renewcommand{\arraystretch}{1.3}
\begin{longtable}{@{}p{3.5cm}p{10cm}@{}}
    \toprule
    \textbf{Abreviatura} & \textbf{Significado} \\
    \midrule
    IoT & Internet de las cosas (\textit{Internet of Things}) \\
    IA  & Inteligencia artificial \\
    OWL & \textit{Web Ontology Language} \\
    RDF & \textit{Resource Description Framework} \\
    % ----- Añade más aquí -----
    \bottomrule
\end{longtable}


\tableofcontents
\listoffigures
\listoftables

% ============================================================
%  CUERPO DEL DOCUMENTO (capítulos)
% ============================================================
\mainmatter

% ============================================================
%  01_introduccion.tex — Capítulo 1: Introducción
% ============================================================

\chapter{Introducción}
\label{cap:introduccion}

% ------------------------------------------------------------------
% TIP: este es el esqueleto. Rellena cada sección con tu contenido.
% ------------------------------------------------------------------

\section{Antecedentes}
\label{sec:antecedentes}

\textit{[Contexto general del problema. ¿Qué se sabe del tema? ¿Qué avances existen
en agricultura de precisión, invernaderos, sensórica e Internet de las cosas?]}

\section{Planteamiento del problema}
\label{sec:planteamiento}

\textit{[Describe el problema concreto que la tesis pretende resolver. ¿Qué dificultad
existe al gestionar un invernadero? ¿Por qué es necesario un enfoque basado en
ontologías / sistemas inteligentes?]}

\section{Justificación}
\label{sec:justificacion}

\textit{[¿Por qué es importante este trabajo? Impacto académico, económico y social.]}

\section{Objetivos}
\label{sec:objetivos}

\subsection{Objetivo general}
\label{subsec:obj-general}

\textit{[Enunciar el objetivo general en una frase.]}

\subsection{Objetivos específicos}
\label{subsec:obj-especificos}

\begin{enumerate}
    \item \textit{[Objetivo específico 1. Ej.: diseñar una ontología del dominio del invernadero.]}
    \item \textit{[Objetivo específico 2.]}
    \item \textit{[Objetivo específico 3.]}
\end{enumerate}

\section{Alcance y limitaciones}
\label{sec:alcance}

\textit{[¿Qué cubre y qué no cubre el trabajo?]}

\section{Estructura de la tesis}
\label{sec:estructura}

\textit{[Breve descripción de cómo está organizado el documento (capítulo por capítulo).]}

% ============================================================
%  02_revision_literatura.tex — Capítulo 2: Revisión de literatura
% ============================================================

\chapter{Revisión de literatura y trabajos relacionados}
\label{cap:estado-del-arte}

\textit{[Este capítulo resume la literatura relacionada. Usa citas con \texttt{\textbackslash cite\{clave\}}
y añade las referencias a \texttt{bibliografia/referencias.bib}.]}

\section{Marco conceptual}
\label{sec:marco-conceptual}

\textit{[Definir los conceptos que sustentan la propuesta: cítricos, invernaderos,
\textit{Tetranychus urticae}, temperatura, humedad relativa, instrumentación e IoT.
Presentar las bases de ontologías y representación del conocimiento que sean
pertinentes para el proyecto.]}

\section{Estado del arte}
\label{sec:estado-del-arte}

\textit{[Describir la búsqueda y selección de literatura sobre monitoreo ambiental,
agricultura de precisión y detección de la plaga. Incorporar únicamente referencias
verificadas en la bibliografía.]}

\subsection{Trabajos relacionados}
\label{subsec:trabajos-relacionados}

\textit{[Comparar los trabajos seleccionados: objetivos, arquitecturas, sensores,
protocolos, métodos de análisis y resultados reportados.]}

\subsection{Brechas existentes}
\label{subsec:brechas}

\textit{[Identificar limitaciones y vacíos sustentados en la revisión de literatura.]}

\subsection{Aportes del proyecto}
\label{subsec:aportes-proyecto}

\textit{[Explicar los aportes propuestos frente a las brechas identificadas, sin
presentarlos como resultados ya demostrados.]}

% Capítulo 3: Necesidades de monitoreo.
\chapter{Necesidades relacionadas con monitoreo de temperatura y humedad relativa
en la detección de \textit{Tetranychus urticae}}
\label{cap:necesidades-monitoreo}

\textit{[Identificar y justificar las necesidades del contexto de estudio, las
variables por observar y los requisitos de monitoreo. Relacionarlos con evidencia
sobre la detección de la plaga; no asumir que las variables ambientales por sí
solas confirman su presencia. Las subsecciones se definirán posteriormente.]}

% Capítulo 4: Arquitectura del sistema.
\chapter{Arquitectura del sistema}
\label{cap:arquitectura}

\textit{[Describir la arquitectura propuesta: componentes de hardware y software,
sensores, adquisición, comunicaciones, almacenamiento, procesamiento y
visualización. Incorporar diagramas y justificar las decisiones de diseño.
Las subsecciones se definirán posteriormente.]}

% Capítulo 5: Plan de pruebas, implementación y evaluación.
\chapter{Plan de pruebas, implementación y evaluación del sistema}
\label{cap:pruebas-implementacion-evaluacion}

\textit{[Documentar el plan de pruebas, los casos y criterios de aceptación,
la implementación de los componentes, la integración de sensores y los
procedimientos y métricas de evaluación. Distinguir las pruebas planificadas
de las efectivamente ejecutadas. Las subsecciones se definirán posteriormente.]}

% Capítulo 6: Análisis de resultados obtenidos.
\chapter{Análisis de resultados obtenidos}
\label{cap:resultados}

\textit{[Presentar e interpretar los resultados obtenidos, contrastarlos con los
objetivos y la literatura, y discutir las limitaciones de la evaluación.
Incorporar tablas y figuras basadas en datos reales. Las subsecciones se
definirán posteriormente.]}

% ============================================================
%  07_conclusiones.tex — Capítulo 7: Conclusiones
% ============================================================

\chapter{Conclusiones}
\label{cap:conclusiones}

\section{Conclusiones generales}
\label{sec:conclusiones-generales}

\textit{[Síntesis de los hallazgos; responde a los objetivos planteados.]}

\section{Aportes del trabajo}
\label{sec:aportes}

\textit{[Contribuciones originales de la tesis.]}

\section{Trabajo futuro}
\label{sec:trabajo-futuro}

\textit{[Líneas de continuación: escalar el sistema, integrar más sensores, optimizar
modelos, etc.]}


% ============================================================
%  PARTE FINAL (destacados, referencias y anexos)
% ============================================================
\chapter*{DESTACADOS}
\markboth{DESTACADOS}{DESTACADOS}
\addcontentsline{toc}{chapter}{DESTACADOS}

\textit{[Sintetizar los aspectos destacados del trabajo una vez se disponga de
resultados. No anticipar hallazgos ni conclusiones sin evidencia.]}


% El encabezado permanece visible aunque todavía no haya citas reales.
\chapter*{REFERENCIAS}
\markboth{REFERENCIAS}{REFERENCIAS}
\addcontentsline{toc}{chapter}{REFERENCIAS}
\printbibliography[heading=none]

% No usar backmatter: desactivaría la numeración alfabética de los anexos.
\chapter*{ANEXOS}
\markboth{ANEXOS}{ANEXOS}
\addcontentsline{toc}{chapter}{ANEXOS}
\appendix
% ============================================================
%  anexo_a.tex — Anexo A: Instrumentos y material complementario
% ============================================================

\chapter{Instrumentos y material complementario}
\label{anexo:a}

\textit{[Material complementario: cuestionarios, tablas de variables, glosario técnico,
especificaciones del hardware, código fuente, etc.]}


\end{document}
